% =============================================================================
% ABSTRACT
% =============================================================================

\begin{abstracts}
    \addcontentsline{toc}{chapter}{ABSTRACT}

    Knowledge graph-constrained decoding eliminates structural hallucination in large language model (LLM) generation by restricting every output token to a verified knowledge graph (KG) path, yet existing frameworks construct their constraint oracles without regard to question semantics or the evolving state of generation.
    %
    This dissertation proposes Dynamic Context-Aware Tries (DCA-Trie), a KG-constrained decoding framework in which the valid token set at each generation step is conditioned jointly on the knowledge graph, the input question, and the partially generated reasoning chain.
    %
    The core limitation addressed is the \emph{permissiveness problem}: leading approaches such as Graph-Constrained Reasoning (GCR) and Decoding on Graphs (DoG) admit all structurally valid KG paths regardless of contextual relevance, forcing the LLM to discriminate among them using its hallucination-prone parametric knowledge.
    %
    DCA-Trie introduces a semantic relevance scorer, implemented with a pretrained sentence transformer, into the constraint oracle to progressively narrow the valid set to paths that are both structurally faithful and semantically aligned with the current reasoning direction.
    %
    Two variants are developed: DCA-Trie v1 applies static semantic filtering at trie construction time, while DCA-Trie v2 performs step-wise dynamic expansion conditioned on the partial generation.
    %
    Both variants are evaluated against GCR and an unconstrained chain-of-thought baseline on the WebQSP and ComplexWebQuestions benchmarks over Freebase, measuring Hits@1, F1, structural faithfulness, and a new Semantic Irrelevance Ratio (SIR) metric introduced in this work.
    %
    The framework is designed to maintain 100\% structural faithfulness while reducing the semantic burden on the LLM at each decoding step, contributing both the DCA-Trie mechanism and the SIR metric to the KG-constrained decoding literature.

\end{abstracts}
